% Supplementary prompting-methods fragment for the LabVerse manuscript.
% Required preamble packages:
% \usepackage{booktabs}
% \usepackage{tabularx}
% \usepackage{array}

\subsection{Prompting schemes and information conditions}
\label{sec:prompting_schemes}

LabVerse evaluates four prompting conditions that separate three design axes:
observation representation (structured versus narrative), state scope (local
versus global), and control regime (interactive versus one-shot). The three
interactive conditions use the same task contract, behavioral instructions,
turn status, action candidates, response schema, and simulator feedback. Their
only intended prompt-level difference is the representation and scope of the
state supplied to the model. The global one-shot planner is a distinct planning
baseline: it receives the initial global state once, emits a complete plan, and
does not observe execution feedback or replan.

\begin{table}[htbp]
\centering
\small
\setlength{\tabcolsep}{4pt}
\renewcommand{\arraystretch}{1.12}
\caption{Primary prompting conditions. ``Global summary'' denotes the compact,
complete whole-laboratory state supplied to the global interactive condition;
``full state'' denotes the complete initial simulator state supplied to the
one-shot planner.}
\label{tab:prompt_conditions}
\begin{tabularx}{\linewidth}{
    >{\raggedright\arraybackslash}p{0.18\linewidth}
    >{\raggedright\arraybackslash}p{0.13\linewidth}
    >{\raggedright\arraybackslash}p{0.24\linewidth}
    >{\raggedright\arraybackslash}p{0.21\linewidth}
    >{\raggedright\arraybackslash}X
}
\toprule
\textbf{Condition} & \textbf{Control} & \textbf{State supplied} &
\textbf{Model output} & \textbf{Post-action information} \\
\midrule
Local structured interactive & One action per call & Structured observation of
the actor's current room and visible entities & One current
\texttt{candidate\_ref}, with free parameters if required & Action result,
recent events, bounded trajectory, and updated local state \\
Local narrative interactive & One action per call & Narrative rendering of the
same local perceptual scope & Same interactive JSON contract & Same feedback
and trajectory contract \\
Global structured interactive & One action per call & Structured local
observation plus compact complete global-state JSON & Same interactive JSON
contract & Same feedback and trajectory contract, plus updated global summary \\
Global one-shot planner & One complete plan per call & Initial local context,
complete initial world-state JSON, and valid-ID inventory & JSON array of
canonical verbs, target IDs, and parameters & None during planning; the frozen
plan is subsequently replayed \\
\bottomrule
\end{tabularx}
\end{table}

\subsubsection{Shared scenario task contract}

For a given scenario, all four conditions receive the same model-facing task
contract. It is generated from the versioned scenario specification and has the
following semantic fields:
\begin{quote}
\begin{verbatim}
Objective: <scenario objective>
Known requirements:
- <requirement 1>
- ...
Required outputs:
- <required output 1>
- ...
Safety constraints:
- <scenario-specific safety constraint 1>
- ...
Intentional unknowns:
- <information that must be discovered through interaction>
- ...
\end{verbatim}
\end{quote}
Only populated sections are rendered. The contract states the operational goal
and any explicitly disclosed requirements; it does not contain completion
ratios, scorer state, criterion weights, hidden oracle state, or instructions
derived from the evaluator. Thus, scenario scoring remains external to the
model prompt.

\subsubsection{Interactive prompt family}

For interactive condition $c$ at decision turn $t$, the model input can be
written as
\begin{equation}
S_c = R \oplus E \oplus D_c \oplus G \oplus F \oplus C,
\qquad
U_c^t = B^t \oplus H^t \oplus T^t \oplus O_c^t
          \oplus W_c^t \oplus A^t,
\label{eq:interactive_prompt}
\end{equation}
where $S_c$ is the system prompt, $U_c^t$ is the current user message,
$\oplus$ denotes ordered concatenation, and the blocks are defined in
Table~\ref{tab:interactive_prompt_variables}. The provider is called with the
system prompt and one current user message; stale full observations and action
lists are not replayed. Prior outcomes are represented only by the bounded
feedback and trajectory blocks.

\begin{table}[htbp]
\centering
\small
\setlength{\tabcolsep}{4pt}
\renewcommand{\arraystretch}{1.10}
\caption{Blocks in the interactive prompt family.}
\label{tab:interactive_prompt_variables}
\begin{tabularx}{\linewidth}{
    >{\centering\arraybackslash}p{0.05\linewidth}
    >{\raggedright\arraybackslash}p{0.20\linewidth}
    >{\raggedright\arraybackslash}p{0.15\linewidth}
    >{\raggedright\arraybackslash}X
}
\toprule
\textbf{Term} & \textbf{Prompt block} & \textbf{Variation} &
\textbf{Contents} \\
\midrule
$R$ & Role & Fixed & Laboratory researcher in a simulated laboratory. \\
$E$ & Environment & Fixed & Turn-based semantics; one action per turn;
candidate list is structural rather than a guarantee of feasibility; failed
actions produce feedback; timed equipment persists across turns. \\
$D_c$ & State-description declaration & By condition & Declares local
structured, local narrative, or local-plus-global structured state access. \\
$G$ & Task contract & By scenario & Objective, known requirements, required
outputs, disclosed safety constraints, and intentional unknowns. \\
$F$ & Response format & Fixed & Exactly one plain JSON object selected by
\texttt{candidate\_ref}; no prose, markup, or wrapper object. \\
$C$ & Action contract & Fixed & Shared safety, PPE, movement, equipment,
waiting, disposal, parameter, and documentation rules. \\
$B^t$ & Turn feedback & By turn & Outcome of the preceding attempted action,
including a blocking or failure reason when applicable. \\
$H^t$ & Compact trajectory & By turn & Bounded summary of previous selected
actions and outcomes; no full conversation replay. \\
$T^t$ & Non-oracle turn status & By turn & Scenario name, used and maximum
steps, remaining steps, and the statement that one selected action consumes one
turn. \\
$O_c^t$ & Current observation & By turn and condition & Current local state in
structured or narrative form. \\
$W_c^t$ & Global-state summary & By turn and condition & Present only for the
global structured interactive condition. \\
$A^t$ & Current action surface & By turn & Exact currently available candidates,
grouped by action type, with targets, fixed/free parameters, hints, durations,
and descriptions. \\
\bottomrule
\end{tabularx}
\end{table}

The following is an abridged, structurally faithful presentation of the
interactive template. Angle-bracketed uppercase names denote runtime
substitutions, not literal model outputs.

\begin{quote}
\begin{verbatim}
SYSTEM
You are a laboratory researcher working in a simulated lab environment.

<environment>
This is a turn-based simulation. On each turn, inspect the state and
the exact current action list, then choose exactly one action. The
candidate list is structural; safety and operational preconditions are
resolved by the simulator. Each selected action consumes one turn.
</environment>

<state_representation>
<CONDITION-SPECIFIC STATE DECLARATION>
</state_representation>

<task_contract>
<SCENARIO TASK CONTRACT>
</task_contract>

<response_format>
Return exactly one plain JSON object matching the current action schema.
The required field is candidate_ref. Include parameters only for free
fields exposed by the selected candidate. Output no prose or wrappers.
</response_format>

<action_contract>
<SHARED SAFETY, GROUNDING, PARAMETER, AND TIMING RULES>
</action_contract>

USER AT TURN t
<turn_feedback><PREVIOUS ACTION OUTCOME></turn_feedback>
<trajectory><BOUNDED ACTION/OUTCOME HISTORY></trajectory>
<turn_status><SCENARIO AND STEP BUDGET></turn_status>
<observation><CURRENT LOCAL OBSERVATION></observation>
<world_state_json><CURRENT GLOBAL SUMMARY></world_state_json>  [global only]
<available_actions><CURRENT CANDIDATE SURFACE></available_actions>
\end{verbatim}
\end{quote}

The condition-specific state block is the sole intended information treatment
within the interactive family:
\begin{itemize}
    \item \textbf{Local structured interactive:} a fielded local observation
          reports the current room, tick, PPE, carried samples, visible
          equipment/materials/samples/personnel, alerts, recent events, and
          visible timed-process state.
    \item \textbf{Local narrative interactive:} the same local perceptual scope
          is rendered as prose. The action surface remains structured and
          identical to the other interactive conditions; narrative rendering
          therefore changes observation representation, not action semantics.
    \item \textbf{Global structured interactive:} the structured local
          observation is augmented by a deterministic compact JSON summary of
          all rooms and adjacency, doors, equipment, materials, samples, and
          personnel. This global summary provides situational awareness, but
          only candidates in the current action surface are executable.
\end{itemize}

At each turn, the action surface supplies a closed set
$A^t=\{a_1^t,\ldots,a_n^t\}$. Each candidate contains a unique turn-local
reference, canonical verb, target when applicable, immutable fixed parameters,
free required and optional parameters, value/type hints, estimated duration,
and a natural-language description. The expected response is
\begin{equation}
Y^t = \{\texttt{candidate\_ref}: r_j^t,
        \ \texttt{parameters}: \theta_j^t\},
\end{equation}
where \texttt{parameters} is omitted when candidate $j$ exposes no free
fields. The candidate reference is the sole branch selector. Optional echoed
fields such as verb or target ID are accepted only when they exactly match the
selected candidate; they are not required and do not select a branch.

Representative output forms are:
\begin{quote}
\begin{verbatim}
{"candidate_ref":"move_1"}

{"candidate_ref":"equip_3",
 "parameters":{"speed_rpm":1000}}
\end{verbatim}
\end{quote}
Fixed parameters are already bound to the candidate and must not be repeated or
modified. This design minimizes redundant generation while guaranteeing that
every interactive action is grounded against the live simulator affordance
surface.

\subsubsection{Global one-shot planning prompt}

The one-shot baseline receives one planning request before execution:
\begin{equation}
S_{\mathrm{OS}} = R_{\mathrm{OS}} \oplus C_{\mathrm{OS}}
  \oplus V \oplus Q \oplus F_{\mathrm{OS}},
\qquad
U_{\mathrm{OS}} = I^0 \oplus W^0 \oplus J^0 \oplus G \oplus L,
\label{eq:oneshot_prompt}
\end{equation}
where $V$ is the canonical action vocabulary, $Q$ contains plan-admissibility
rules, $I^0$ is the initial local context and turn-zero affordances, $W^0$ is
the complete initial world-state JSON, $J^0$ is the valid-ID and parameter-name
inventory, and $L$ is the run step budget. The task contract $G$ is unchanged.

\begin{quote}
\begin{verbatim}
SYSTEM
You are a laboratory researcher planning a series of actions.
Given the complete lab configuration, task, action vocabulary, and
valid-ID inventory, output a COMPLETE plan as a JSON array.

<safety_rules><SHARED LAB SAFETY AND TIMING RULES></safety_rules>
<action_vocabulary><CANONICAL VERBS AND SIGNATURES></action_vocabulary>
<plan_admissibility_rules>
The plan must not be empty or passive filler. Use wait only to advance
a running timed process.
</plan_admissibility_rules>
Use only IDs and parameter names from the supplied inventory.
Do not emit candidate_ref values. Output only the JSON array.

USER
<initial_context><TURN-ZERO LOCAL STATE AND ACTIONS></initial_context>
<full_lab_state_json><COMPLETE INITIAL WORLD STATE></full_lab_state_json>
<valid_id_inventory_json><VALID IDS AND PARAMETERS></valid_id_inventory_json>
<task_contract><SCENARIO TASK CONTRACT></task_contract>
<run_constraints><MAXIMUM STEPS; ONE ACTION PER STEP></run_constraints>
\end{verbatim}
\end{quote}

The output is a sequence
\begin{equation}
Y_{\mathrm{OS}} =
\left[(v_1,z_1,\theta_1),\ldots,(v_k,z_k,\theta_k)\right],
\end{equation}
serialized as JSON objects with \texttt{verb}, optional
\texttt{target\_id}, and optional \texttt{parameters} fields. The canonical
vocabulary comprises movement; door operations; PPE changes; equipment start
and stop; pipette transfer; sample retrieval and storage; material disposal;
waiting; and observation logging. Legacy aliases are explicitly disallowed.

Candidate references are deliberately absent from one-shot output. They are
turn-local labels generated from the current state, and future labels cannot be
known before earlier planned actions have changed that state. During replay,
each planned verb/target/parameter tuple must bind to the live action surface at
that step. Failure to bind is recorded as a grounding failure. Consequently,
the most direct control-regime comparison is between global structured
interactive control and global one-shot planning; the two local interactive
conditions additionally test partial observability.

\subsubsection{Multi-agent prompt extension}

The shared-equipment scenario uses the same four information conditions with a
team coordinator. Interactive prompts contain one local observation per
technician and one grouped team action surface. Each candidate reference is
namespaced by technician and uniquely selects both the actor and action; the
coordinator still returns exactly one action per turn:
\begin{quote}
\begin{verbatim}
{"candidate_ref":"t1_equip_2"}
\end{verbatim}
\end{quote}
The global interactive variant additionally receives the whole-laboratory
summary. For one-shot team planning, every plan element adds the required
\texttt{agent\_id} field:
\begin{quote}
\begin{verbatim}
[{"agent_id":"tech_1","verb":"move_to","target_id":"wet_lab"},
 {"agent_id":"tech_2","verb":"don_ppe",
  "parameters":{"ppe_type":"lab_coat"}}]
\end{verbatim}
\end{quote}
The team contract additionally requires safe scheduling, avoidance of equipment
contention, and exactly one centrifuge run per technician. All team actions are
executed as ordered transitions on one shared world state and event timeline.

\subsubsection{Retries, parsing, and prompt audit}

Malformed or protocol-invalid model output may be retried at most twice after
the initial call. A retry includes the parsing or grounding error and restates
the relevant JSON contract; it does not advance simulator time, submit an
action, or consume an environment step. The same maximum of two corrective
retries applies to interactive decisions and one-shot plan generation. Output
normalization and JSON repair are parser-side robustness mechanisms and do not
alter the state information or candidate set shown to the model. Retry counts,
normalization events, provider errors, and the raw model responses are retained
in the run artifacts.

To guard experimental parity, the benchmark audit verifies that the three
interactive conditions share the same task contract, environment instructions,
response format, action contract, available-action text, turn status, feedback,
trajectory, and dynamic action schema for a matched state. It also checks that
local prompts do not contain global-state blocks and that prompt-visible status
does not expose scorer-derived progress. Prompt snapshots and prompt-family
hashes are stored with experiment artifacts, allowing the exact evaluated
prompt surface to be inspected and prompt drift to be detected.
